\section{Reproducibility}\label{app:repro}

The code, frozen evidence and report sources are in the public repository
\url{https://github.com/brmurrell3/macro-placer}. Commands below run from the
repository root with the project's virtual environment.

\paragraph{Source revisions.} Three revisions anchor the work. The competition
snapshot \texttt{2ab92a33d352} (22 May 2026) is the last competition-era commit; its
Python source closure is frozen as \path{writeup/data/publication/competition-source.tar},
described by \path{source_snapshot.json} with the archive digest and file list. The
upstream challenge revision \texttt{996193c83eaa} (14 August 2026) carries the final
leaderboard. The TILOS MacroPlacement submodule is pinned at \texttt{45a721d01dfe};
six canonical evaluator, loader and validation files are byte-identical between the
local base and the final challenge revision. The ranked single pipeline is most
likely commit \texttt{5485310}; Appendix~\ref{app:records} gives the basis for that
inference.

\paragraph{Package extraction.} The competition entries imported code from the
experiment tree, so copying a submission directory onto a clean checkout of the
challenge was not enough to run it. The post-competition cleanup extracts the active
dependency closure into the package \path{thinkorplace/} with three parts,
\path{thinkorplace.common}, \path{thinkorplace.v1} and \path{thinkorplace.v2}, using
relative imports; the submission entry points remain as thin adapters. Runtime
imports no longer touch \path{experiments/} or \path{writeup/}.
\path{thinkorplace/MIGRATION.json} records the origin and source digest of every
module. The package keeps the optimization formulas and adds the repairs of
Appendix~\ref{app:verification}; a test checks that the extracted primitives
preserve the original algorithms' syntax trees.

\paragraph{Frozen evidence and generators.} \path{writeup/data/publication/} holds
the evidence, and its \path{README.md} states the scope of each record:
\begin{itemize}
  \item \path{raw/}: the three historical suite records, the 34 curvature-stage
        records and the 34 saved placement tensors, with SHA-256 values in
        \path{historical_summary.json};
  \item the evaluator microbenchmarks, curvature diagnostic and follow-up,
        rescoring and geometry records, topology diagnostics, tail-bound and
        error-margin arithmetic, and the frozen leaderboard;
  \item \path{ibm_repair_equivalence.json}, written by
        \path{scripts/measure_publication_ibm_repair_equivalence.py}: on all 17 IBM
        inputs, the 20 May extraction code (commit \texttt{5485310}) and the
        corrected runtime agree on pin ownership, smooth net tensors and initializer
        edges, and the record counts the input soft clusters that the initial clamp
        moves; no optimization or scoring is run;
  \item summaries of the three engineering batches, whose primary run records
        (manifests, requests, results, logs, placements and checkpoint hashes) are
        under \path{writeup/runs/}, as are those of the stopped first reproduction
        attempt (\path{writeup/runs/ranked-reproduction-20260930/});
  \item \path{corrected-runtime-*.tar.gz}: deterministic archives of each
        post-competition runtime, with per-file digests in the matching
        \path{corrected_runtime_snapshot_*.json}.
\end{itemize}
Generators read these records and write only derived summaries, tables and figures;
the historical freezer refuses to replace a frozen raw record with different bytes.
No generator runs a solver, and the scripts that produced the diagnostics are separate
from the generators.

\paragraph{Runner.} The study runner (\path{scripts/run_publication_study.py} with
\path{scripts/publication_worker.py}) implements the following workflow:
\begin{enumerate}
  \item create an immutable run manifest before any solver is constructed;
  \item record input and initialization hashes and the explicit device choice;
  \item save canonical incumbents atomically at every meaningful phase;
  \item propagate one monotonic solver deadline and enforce an external
        process-group timeout;
  \item validate the final candidate independently and retain every failure;
  \item preserve final coordinates and phase metadata for deterministic
        post-processing.
\end{enumerate}
A run's identity covers executable source, dependencies, inputs and constructor
settings, so editing the paper does not invalidate a run while any runtime change
does. Source is hashed again after execution. Validity (shape, finite values,
containment, fixed objects, zero hard overlaps), study eligibility (return within the
external limit with matching source and environment) and pipeline completion are
separate fields: a valid recovered checkpoint can be an eligible anytime result
without a completed pipeline.

\paragraph{Tests.} Runtime archive \texttt{5aa4c719cf18} passed 64 focused tests in
five modules; the receipt \path{verification/final-tests-20260906.json} binds the
archive to the test log:
\begin{itemize}
  \item \path{test_submission_safety.py} (30): deadline propagation, zero join
        timeouts after expiry, child termination and cleanup, rejection of invalid
        lanes, checkpoint retention and atomic replacement, phase-outcome
        reporting, hierarchical pin ownership in four code copies, and budget guards
        for the curvature stages and the eigensolver;
  \item \path{test_thinkorplace_package.py} (7): imports against a pristine
        upstream without the research tree, legalizer clamping of soft clusters,
        parity of runtime validation with the canonical check on float32 endpoints
        and randomized geometry, and preservation of the extracted algorithms;
  \item \path{test_postpilot_repairs.py} (12): soft-cluster bounds, fixed obstacles,
        float32 export, infeasible fixed references, unresolved hard geometry, NG45
        directory aliases, and retained diagnostics for rejected initializers and
        fatal results;
  \item \path{test_sdf_topology.py} (4): initializer edge extraction for flat and
        hierarchical netlists in both code copies;
  \item \path{test_publication_runner.py} (11): explicit parameters, worker
        validation, checkpoint-recovery semantics, process-group timeouts, resume
        identity, dry runs and checkpoint inventories.
\end{itemize}
These are engineering checks of execution and validation, not measurements of
placement quality.

\paragraph{Commands.} Regenerate the tables and figures from the frozen records:
\begin{verbatim}
.venv/bin/python scripts/freeze_publication_evidence.py
.venv/bin/python scripts/render_publication_measurements.py
.venv/bin/python scripts/render_publication_curvature.py
.venv/bin/python scripts/render_publication_pilot.py
.venv/bin/python scripts/render_publication_pilot.py --profile repair
.venv/bin/python scripts/render_publication_integration.py
.venv/bin/python scripts/render_publication_devrecords.py
.venv/bin/python scripts/render_publication_v1.py
\end{verbatim}
Build and check the report:
\begin{verbatim}
make -C writeup/latex
.venv/bin/python scripts/check_publication_paper.py
\end{verbatim}
Run the focused tests:
\begin{verbatim}
.venv/bin/python -m pytest test/test_submission_safety.py \
  test/test_thinkorplace_package.py test/test_postpilot_repairs.py \
  test/test_sdf_topology.py test/test_publication_runner.py
\end{verbatim}
Run a study from a configuration under \path{writeup/study/}, for example the pilot,
or \path{writeup/study/ranked_reproduction.json} for the ranked configuration:
\begin{verbatim}
.venv/bin/python scripts/run_publication_study.py \
  --config writeup/study/pilot.json
\end{verbatim}
Benchmark inputs and the evaluator come from the pinned submodule; every run records
their hashes.
